\documentclass[11pt,a4paper]{article}
\usepackage[utf8]{inputenc}
\usepackage[T1]{fontenc}
\usepackage{lmodern}
\usepackage{amsmath,amssymb,amsthm}
\usepackage[margin=2.5cm]{geometry}
\usepackage{xcolor}
\usepackage{hyperref}
\usepackage{booktabs}

\newcommand{\PiTT}{\Pi_{\mathrm{TT}}}
\newcommand{\Fhat}{\hat{F}}
\newcommand{\dd}{\mathrm{d}}

\newtheorem{theorem}{Theorem}
\newtheorem{corollary}[theorem]{Corollary}
\newtheorem{lemma}[theorem]{Lemma}
\newtheorem{remark}[theorem]{Remark}

\hypersetup{
  colorlinks=true,
  linkcolor=blue!60!black,
  pdftitle={Ghost Inevitability in the Spectral Action},
  pdfauthor={David Alfyorov},
}

\title{%
  \textbf{Ghost Inevitability Theorem}\\[4pt]
  \large Structural Ghost Poles in the Spectral Action
  with Standard Model Matter Content
}
\author{David Alfyorov}
\date{March 2026}

\begin{document}
\maketitle

\begin{center}
\small\textit{Credits: Aliaksandr Samatyia contributed research-assistance
and workflow support.}
\end{center}

\begin{abstract}
We prove that the one-loop graviton propagator in Spectral Causal Theory
(SCT) necessarily contains ghost poles for any spectral action
$S = \mathrm{Tr}\,f(D^2/\Lambda^2)$ with Standard Model matter content.
The result follows from three facts:
(i)~the master function $\varphi(z)$ and the normalized form factor
$\Fhat_1(z)$ depend only on the heat-kernel structure, not on the
choice of~$f$;
(ii)~the Weyl${}^2$ coefficient $\alpha_C = 13/120$ is positive and
determined solely by SM multiplicities;
(iii)~the intermediate value theorem applied to the entire function
$\PiTT(z)$, which satisfies $\PiTT(0) = 1 > 0$ while
$\PiTT(z) \to -83/6 < 0$ as $z \to +\infty$ and
$\PiTT(z) \to -\infty$ as $z \to -\infty$.
No modification of~$f$ can eliminate the ghost poles.  The only escape
is to change the matter content such that
$N_s + 12\,N_v \leq 6\,N_D$, which fails for the SM by a margin of~13
degrees of freedom.
\end{abstract}

\section{Setup}

The one-loop graviton propagator denominator in SCT is
\begin{equation}\label{eq:PiTT}
  \PiTT(z) = 1 + \frac{13}{60}\,z\,\Fhat_1(z),
\end{equation}
where $z = \Box/\Lambda^2$ is the dimensionless d'Alembertian and
$\Fhat_1(z) = F_1(z)/F_1(0)$ is the normalized spin-2 form factor.

The form factor $F_1(z)$ is constructed from the heat-kernel master
function
\begin{equation}\label{eq:phi}
  \varphi(z)
  = \int_0^1 e^{-\xi(1-\xi)z}\,\dd\xi
  = e^{-z/4}\sqrt{\frac{\pi}{z}}\,\mathrm{erfi}\!\left(
    \frac{\sqrt{z}}{2}\right),
\end{equation}
which is an entire function of~$z$ with $\varphi(0) = 1$.

The coefficient $13/60 = 2\,\alpha_C$, where
\begin{equation}\label{eq:alphaC}
  \alpha_C
  = N_s\,\beta_W^{(0)}
    + N_D\,\beta_W^{(1/2)}
    + N_v\,\beta_W^{(1)}
  = \frac{N_s}{120} - \frac{N_D}{20} + \frac{N_v}{10},
\end{equation}
with the SM multiplicities $N_s = 4$, $N_D = 45/2$, $N_v = 12$,
giving $\alpha_C = 13/120$.

\begin{remark}[$f$-independence]
The function $\varphi(z)$ in~\eqref{eq:phi} depends only on the
eigenvalue distribution of $D^2$ through the heat kernel
$K(t) = \mathrm{Tr}\,e^{-tD^2}$, not on the spectral function
$f(x)$ in $S = \mathrm{Tr}\,f(D^2/\Lambda^2)$.
The per-spin form factors $h_C^{(s)}(z)$ and $h_R^{(s)}(z)$ are
algebraic functions of $\varphi(z)$ and~$z$.
The choice of $f$ affects only the overall normalization $F_1(0)$,
which cancels in the ratio $\Fhat_1(z) = F_1(z)/F_1(0)$.
Therefore, $\PiTT(z)$ is entirely determined by the heat-kernel
structure and the SM particle content, independently of~$f$.
\end{remark}


\section{Ghost Inevitability Theorem}

\begin{theorem}[Ghost inevitability]
\label{thm:main}
For any spectral function $f(x)$ and any matter content with
$\alpha_C > 0$ (equivalently, $N_s + 12\,N_v > 6\,N_D$),
the propagator denominator $\PiTT(z)$ has at least two real zeros:
one at some $z_0 > 0$ (Euclidean/spacelike ghost) and one at some
$z_L < 0$ (Lorentzian/timelike ghost).
\end{theorem}

\begin{proof}
We establish three properties of $\PiTT(z)$:

\textbf{(i)~$\PiTT(0) = 1 > 0$.}
At $z = 0$: $\Fhat_1(0) = 1$ by normalization, so
$\PiTT(0) = 1 + 0 = 1$.

\textbf{(ii)~$\PiTT(z) \to -83/6 < 0$ as $z \to +\infty$.}
From the verified Phase~3 UV asymptotic
$z\,\alpha_C(z) \to -89/12$ as $z \to +\infty$
(where $\alpha_C(z) = \sum_s N_s\,h_C^{(s)}(z)$ is the
$z$-dependent Weyl${}^2$ coefficient), we obtain
\begin{equation}
  \PiTT(z) = 1 + 2\,z\,\alpha_C(z)
  \;\xrightarrow{z \to +\infty}\;
  1 - \frac{89}{6} = -\frac{83}{6} \approx -13.83.
\end{equation}
Numerically verified: $\PiTT(10^4) = -13.833$ to 4~significant
figures.

\textbf{(iii)~$\PiTT(z) \to -\infty$ as $z \to -\infty$.}
For $z < 0$ (the negative real axis), $\varphi(z) = \int_0^1 e^{|z|\xi(1-\xi)}\,\dd\xi$
grows as $\varphi(z) \sim e^{|z|/4}\sqrt{\pi/|z|}$ (dominated by
the saddle point at $\xi = 1/2$).  This is distinct from the
positive-real-axis UV limit, where $\varphi(z) \sim 2/z$ as $z \to +\infty$.
The form factors $h_C^{(s)}(z)$ therefore grow exponentially only on
the negative real axis, so $F_1(z)$ grows exponentially for
$z \ll 0$.  Since $z < 0$ and $F_1(z) > 0$ for $z < 0$, the product
$z \cdot \Fhat_1(z) \to -\infty$, and therefore
$\PiTT(z) \to -\infty$.
Numerically: $\PiTT(-10) = -182$, $\PiTT(-50) \approx -1.7 \times 10^7$.

\medskip
Since $\PiTT(z)$ is an entire function (hence continuous on~$\mathbb{R}$)
satisfying
\[
  \PiTT(0) = 1 > 0,
  \qquad
  \lim_{z \to +\infty} \PiTT(z) = -\frac{83}{6} < 0,
  \qquad
  \lim_{z \to -\infty} \PiTT(z) = -\infty,
\]
the intermediate value theorem guarantees:
\begin{itemize}
  \item $\exists\, z_0 \in (0, +\infty)$ with $\PiTT(z_0) = 0$
    \quad (Euclidean ghost, $k^2 < 0$, spacelike),
  \item $\exists\, z_L \in (-\infty, 0)$ with $\PiTT(z_L) = 0$
    \quad (Lorentzian ghost, $k^2 > 0$, timelike).
\end{itemize}
Neither zero can be removed by modifying~$f$, since $\PiTT(z)$ is
$f$-independent.
\end{proof}


\begin{corollary}[SM ghost margin]
\label{cor:margin}
For the Standard Model, $N_s + 12\,N_v = 148$ and $6\,N_D = 135$.
The ghost-inevitability condition $N_s + 12\,N_v > 6\,N_D$ is
satisfied with a margin of~$13$ degrees of freedom.  To achieve
$\alpha_C \leq 0$ would require adding at least $\lceil 13/6 \rceil + 1
= 3$ Dirac fermion doublets ($\approx 5$ Weyl fermions) beyond the SM
content.
\end{corollary}


\section{Matter Content Analysis}

\begin{center}
\begin{tabular}{lcccl}
\toprule
Theory & $N_s + 12\,N_v$ & $6\,N_D$ & $\alpha_C$ & Ghost? \\
\midrule
Standard Model & 148 & 135 & $13/120 > 0$ & \textbf{Yes} \\
Pure scalar ($N_s > 0$) & $N_s$ & 0 & $> 0$ & Yes \\
Pure vector ($N_v > 0$) & $12\,N_v$ & 0 & $> 0$ & Yes \\
Pure fermion ($N_D > 0$) & 0 & $6\,N_D$ & $< 0$ & \textbf{No} \\
SM + 3 Dirac doublets & 148 & 153 & $-5/120 < 0$ & No \\
\bottomrule
\end{tabular}
\end{center}

\begin{remark}[Physical interpretation]
Fermions contribute \emph{negatively} to $\alpha_C$ because the
one-loop fermion determinant carries an overall minus sign (closed
fermion loop).  A theory with sufficiently many fermions relative to
bosons can have $\alpha_C \leq 0$, in which case
$\PiTT(z \to +\infty) \geq 1$ and the positive-axis ghost is absent.
Whether a negative-axis ghost survives in this case depends on the
detailed structure of $\Fhat_1(z)$ and requires separate analysis.
\end{remark}


\section{Numerical Verification}

\begin{center}
\begin{tabular}{rcc}
\toprule
$z$ & $\PiTT(z)$ & Sign \\
\midrule
$-50$ & $-1.72 \times 10^7$ & $-$ \\
$-10$ & $-182.4$ & $-$ \\
$-5$ & $-19.64$ & $-$ \\
$-2$ & $-1.395$ & $-$ \\
$-1.281$ & $-4.3 \times 10^{-4}$ & $\approx 0$ \\
$-1$ & $+0.361$ & $+$ \\
$0$ & $+1.000$ & $+$ \\
$+0.5$ & $+1.023$ & $+$ \\
$+2$ & $+0.328$ & $+$ \\
$+2.415$ & $-1.4 \times 10^{-4}$ & $\approx 0$ \\
$+3$ & $-0.524$ & $-$ \\
$+10$ & $-7.314$ & $-$ \\
$+10^4$ & $-13.833$ & $-$ \\
\bottomrule
\end{tabular}
\end{center}

The zeros are located at $z_L = -1.28070\ldots$ (Lorentzian ghost)
and $z_0 = +2.41484\ldots$ (Euclidean ghost), matching the MR-2 ghost
catalogue to 15+ significant figures.


\section{Implications}

\begin{enumerate}
  \item \textbf{No escape via $f$:} The ghost poles are a structural
    consequence of the spectral action with SM matter, not an artifact
    of a particular choice of spectral function.

  \item \textbf{The ghost problem is real:} Any resolution must come
    from the quantum treatment of the ghost degrees of freedom (fakeon,
    Donoghue--Menezes, PT symmetry), not from modifying the classical
    action.

  \item \textbf{Predictivity:} The inevitability of ghosts is the
    flip side of SCT's high predictivity.  The form factors are
    \emph{derived}, not \emph{postulated}.  A theory that postulates
    ghost-free form factors (e.g., infinite-derivative gravity with
    $e^{-\Box/\Lambda^2}$) avoids ghosts but loses predictivity.

  \item \textbf{BSM sensitivity:} The ghost margin ($148 - 135 = 13$)
    is remarkably small.  Adding $\sim 5$ Weyl fermions to the SM
    would eliminate $\alpha_C > 0$.  This makes the ghost problem
    sensitive to BSM physics near the Planck scale.
\end{enumerate}

\end{document}
